\subsection*{Week of 2026-11-30: counting and configurations}

Second rail: invariants, games and hidden symmetry. The second timed half paper is on Wednesday 2026-12-02. Thursday and Friday: no new problems; read the Misses and Known sections of \texttt{readout.md} once each, and the Weaknesses section of \texttt{profile.md}.

\begin{center}\small
\begin{tabular}{@{}c >{\raggedright\arraybackslash}p{3.1cm} >{\raggedright\arraybackslash}p{4.8cm} >{\raggedright\arraybackslash}p{4.8cm}@{}}
\toprule
Rung & Problem & Object and ask & Move \\
\midrule
1 & \seminar{PS 10 \#1, \#4} and \seminar{\#7} & Three exercise-like configuration problems, as the set itself says. & Colourings and sums. \\
2 & \putnam{2018 B1} & A set of lattice vectors; remove one so the rest split into two equal-sum halves. & Hidden symmetry: pair things up. \\
3 & \putnam{2019 B1} & Lattice points on circles; count squares. & Count by a bijection with something you can list. \\
4 & \putnam{2023 B1} & Coins sliding on a grid; which configurations are reachable. & An invariant of the moves. \\
5 & \putnam{2024 B1} & A grid with a prescribed set of values along a transversal; feasibility. & Invariants again, with a construction for the other direction. \\
6 & \putnam{2020 B1} & A signed sum over integers by the parity of their binary digits; a residue. & Tagged number theory; the hidden symmetry of the sign pattern carries it. Cite the rung where you first paired terms off. \\
\midrule
opt & \putnam{2017 A4} & Quiz scores; split a class into two halves with equal averages. & \\
opt & \putnam{2020 B2} & A peg-moving game; when the first player wins. & \\
\bottomrule
\end{tabular}
\end{center}

Seminar: \seminar{PS 10 \#3, \#5, \#16}; \seminar{PS 1 \#1, \#2, \#3}.

\textbf{Mock 2, Wednesday 2026-12-02.} The 2014 B paper cut the same way: B1, B3 and B5, a fifteen-minute break, then B2, B4 and B6. Same protocol.
